\chapter{AI Explainability Standard}\label{record-033}

\noindent{\footnotesize\color{BookMuted}\path{docs/strategy/ai-explainability-standard.md}\par}\smallskip

\emph{Design note for explicit, governed AI-mediated UX in Bijou surfaces}

\section{Why this exists}\label{record-033-why-this-exists}

If Bijou is going to support AI-mediated interfaces, suggestions, assistants, or explainability surfaces, it needs a product standard that is stricter than ``show some generated text.''

The \texttt{x\allowbreak{}y\allowbreak{}p\allowbreak{}h} TUI already demonstrates a strong pattern:

\begin{itemize}
\tightlist
\item
  visible \texttt{[\allowbreak{}A\allowbreak{}I\allowbreak{}]} labeling
\item
  explicit explainability sections
\item
  rationale and evidence
\item
  expected next action
\item
  confidence-aware posture
\end{itemize}

Bijou should adopt that seriousness as doctrine.

AI output should never feel like:

\begin{itemize}
\tightlist
\item
  hidden automation
\item
  unexplained authority
\item
  decorative speculation
\end{itemize}

\section{Product stance}\label{record-033-product-stance}

AI-mediated UI in Bijou should be:

\begin{itemize}
\tightlist
\item
  explicit
\item
  attributable
\item
  explainable
\item
  actionable
\item
  governed
\end{itemize}

AI should surface as advisory or assistive unless the app explicitly states otherwise.

Users should be able to answer:

\begin{itemize}
\tightlist
\item
  what is AI-generated here?
\item
  why is it being shown?
\item
  what evidence or basis does it rely on?
\item
  what am I expected to do with it?
\item
  how confident is it?
\end{itemize}

\section{Core principles}\label{record-033-core-principles}

\subsection{1. Mark AI visibly}\label{record-033-mark-ai-visibly}

AI-mediated content should use an explicit \texttt{[\allowbreak{}A\allowbreak{}I\allowbreak{}]} marker or equivalent canonical treatment.

The user should not have to infer that a suggestion, explanation, or summary came from an agent.

That marker should be part of a canonical explainability component/pattern, not a one-off badge.

DL-005 establishes that first concrete seam through \texttt{e\allowbreak{}x\allowbreak{}p\allowbreak{}l\allowbreak{}a\allowbreak{}i\allowbreak{}n\allowbreak{}a\allowbreak{}b\allowbreak{}i\allowbreak{}l\allowbreak{}i\allowbreak{}t\allowbreak{}y\allowbreak{}(\allowbreak{})}.

\subsection{2. Explain why}\label{record-033-explain-why}

AI output should include or make available:

\begin{itemize}
\tightlist
\item
  why it was suggested
\item
  what it is based on
\item
  what decision or workflow it is trying to support
\end{itemize}

\subsection{3. Separate evidence from recommendation}\label{record-033-separate-evidence-from-recommendation}

Recommendation should not masquerade as fact.

Bijou should encourage surfaces that distinguish:

\begin{itemize}
\tightlist
\item
  observation
\item
  evidence
\item
  recommendation
\item
  next action
\end{itemize}

\subsection{4. Confidence must be interpretable}\label{record-033-confidence-must-be-interpretable}

If confidence is shown, it should be attached to explainable reasoning rather than used as mystical authority.

Confidence without context increases user anxiety and misuse.

\subsection{5. AI should lead to governed next actions}\label{record-033-ai-should-lead-to-governed-next-actions}

A useful AI surface should make it clear what comes next:

\begin{itemize}
\tightlist
\item
  review
\item
  accept
\item
  reject
\item
  inspect more evidence
\item
  do nothing
\end{itemize}

\subsection{6. AI should reduce ambiguity, not add it}\label{record-033-ai-should-reduce-ambiguity-not-add-it}

An explainability surface is successful when it lowers uncertainty.

It is unsuccessful when it adds more jargon or produces an attractive wall of text without operational value.

\section{Canonical explainability fields}\label{record-033-canonical-explainability-fields}

Bijou should gradually standardize around explainability structures like:

\begin{itemize}
\tightlist
\item
  label: \texttt{[\allowbreak{}A\allowbreak{}I\allowbreak{}]}
\item
  artifact kind
\item
  source
\item
  source mode
\item
  rationale / why
\item
  evidence basis
\item
  confidence
\item
  expected next action
\item
  governance or review note
\end{itemize}

This does not have to be one monolithic component immediately.

It does need to become a consistent standard.

\section{DOGFOOD and docs implication}\label{record-033-dogfood-and-docs-implication}

DOGFOOD should eventually teach AI-mediated surfaces the same way it teaches other component families:

\begin{itemize}
\tightlist
\item
  labeled clearly
\item
  explained honestly
\item
  shown with rationale and next-action posture
\end{itemize}

This is especially important if Bijou later grows:

\begin{itemize}
\tightlist
\item
  AI-assisted blocks
\item
  AI-authored settings or recommendations
\item
  design/debug assistants
\end{itemize}

\section{Relationship to XYPH}\label{record-033-relationship-to-xyph}

This direction is influenced by the stronger explainability posture already present in \texttt{x\allowbreak{}y\allowbreak{}p\allowbreak{}h}, including:

\begin{itemize}
\tightlist
\item
  explicit \texttt{[\allowbreak{}A\allowbreak{}I\allowbreak{}]} labeling in TUI views
\item
  structured explainability bodies
\item
  rationale and evidence sections
\item
  confidence-aware governance posture
\end{itemize}

Bijou should learn from that standard rather than reinventing a weaker one.

\section{Relationship to other doctrine}\label{record-033-relationship-to-other-doctrine}

This note should be read alongside:

\begin{itemize}
\tightlist
\item
  \href{https://github.com/flyingrobots/bijou/blob/cbd70612140e7a450d764b3c7be1c6af878b99bc/docs/strategy/bijou-ux-doctrine.md}{Bijou UX Doctrine}
\item
  \href{https://github.com/flyingrobots/bijou/blob/cbd70612140e7a450d764b3c7be1c6af878b99bc/docs/strategy/content-guide.md}{Content Guide}
\item
  \href{https://github.com/flyingrobots/bijou/blob/cbd70612140e7a450d764b3c7be1c6af878b99bc/docs/strategy/accessibility-and-assistive-modes.md}{Accessibility and Assistive Modes}
\end{itemize}

\section{Acceptance bar}\label{record-033-acceptance-bar}

This direction is working when:

\begin{itemize}
\tightlist
\item
  AI-mediated surfaces in Bijou are visibly marked
\item
  explainability becomes a product standard rather than ad hoc copy
\item
  users can understand why an AI artifact is present and what they should do next
\item
  the framework encourages governed, attributable AI experiences instead of mysterious ones
\end{itemize}

